% from ega4-8-fr.tex
\oldpage[IV-3]{37}
restriction de $f_n$ à $f_n^{-1}(Y')$~; mais pour tout ouvert $V$ de $X_n$,
on a $f_n(V\cap f_n^{-1}(Y'))=f_n(V)\cap Y'$, et comme par hypothèse $f_n(V)$
est ouvert dans $Y_n$, cela montre que $f'$ est aussi un morphisme ouvert.
Lorsque $A'$ est quelconque, il peut être considéré comme limite inductive de
ses sous-$A$-algèbres de type fini $A'_\lambda$, et le fait que $f'$ soit un
morphisme ouvert résulte de ce qui précède et de (8.10.1).
\end{proof}

\begin{proposition}[8.10.3]
\EGAReaderTarget{ega4.c04.prop.8-10-3}{IV.8.10.3-fr}
Supposons qu'il existe $\alpha$ tel que~: 1\textsuperscript{o} $S_\alpha$ soit
quasi-compact~; 2\textsuperscript{o} les morphismes $X_\alpha\to S_\alpha$,
$Y_\alpha\to S_\alpha$ soient quasi-compacts et le morphisme
$Y_\alpha\to S_\alpha$ quasi-séparé~; 3\textsuperscript{o} pour
$\alpha\leq\lambda\leq\mu$, les morphismes
$u_{\lambda\mu}:S_\mu\to S_\lambda$ soient plats~; 4\textsuperscript{o}
$\overline{f_\alpha(X_\alpha)}$ soit constructible dans $Y_\alpha$. Alors,
pour que $f$ soit dominant, il faut et il suffit qu'il existe
$\lambda\geq\alpha$ tel que $f_\lambda$ soit dominant.
\end{proposition}

\begin{proof}
Les hypothèses entraînent que $Y_\alpha$ est quasi-compact et que le morphisme
$f_\alpha$ est quasi-compact (I, 1.2.4)~; par suite
$\overline{f_\alpha(X_\alpha)}=Z_\alpha$ est pro-constructible (I, 1.9.5,
(vii)) dans $Y_\alpha$. Si l'on pose
$Z_\lambda=w_{\alpha\lambda}^{-1}(Z_\alpha)$ pour $\lambda\geq\alpha$ et
$Z=w_\alpha^{-1}(Z_\alpha)$, on a donc
$Z_\lambda=\overline{f_\lambda(X_\lambda)}$ et
$Z=\overline{f(X)}$ (I, 3.4.8), et $Z_\lambda$ est pro-constructible dans
$Y_\lambda$ (I, 1.9.5, (vii)). Il suffit alors d'appliquer (8.3.13) en
remplaçant $S_\lambda$, $Z'_\lambda$ et $Z''_\lambda$ par $Y_\lambda$,
$Y_\lambda$ et $Z_\lambda$ respectivement.
\end{proof}

\begin{proposition}[8.10.4]
\label{IV.8.10.4-fr}
Supposons qu'il existe $\alpha$ tel que $Y_\alpha$ soit quasi-compact et
$f_\alpha$ de type fini et quasi-séparé. Pour que le morphisme $f$ soit
séparé, il faut et il suffit qu'il existe $\lambda\geq\alpha$ tel que
$f_\lambda$ soit séparé.
\end{proposition}

\begin{proof}
La question étant locale sur $Y_\alpha$ (puisque $Y_\alpha$ est quasi-compact
et $L$ filtrant), on peut se borner au cas où $Y_\alpha$ est affine, donc
quasi-séparé, et l'hypothèse entraîne que $X_\alpha$ (donc les $X_\lambda$ et
$X$) sont quasi-compacts et quasi-séparés. Posons
$X'_\lambda=X_\lambda\times_{Y_\lambda}X_\lambda$ pour
$\lambda\geq\alpha$ et $X'=X\times_YX$~; on a
$X'_\lambda=X'_\alpha\times_{Y_\alpha}Y_\lambda$ et
$X'=X'_\alpha\times_{Y_\alpha}Y$~; le morphisme première projection
$X'_\alpha\to X_\alpha$ est quasi-compact et quasi-séparé par hypothèse
(I, 1.2.3, (iii)), donc $X'_\lambda$ est quasi-compact et quasi-séparé.
Notons maintenant que si l'on désigne par $\Delta_\lambda$ (resp. $\Delta$)
la diagonale de $X_\lambda\times_{Y_\lambda}X_\lambda$ (resp. de
$X\times_YX$), il résulte de (I, 5.3.4 et 3.4.8) que $\Delta_\mu$ (resp.
$\Delta$) est l'image réciproque de $\Delta_\lambda$ par l'application
$v'_{\lambda\mu}:X'_\mu\to X'_\lambda$ (resp.
$v'_\lambda:X'\to X'_\lambda$). D'autre part $\Delta_\lambda$ est
constructible dans $X'_\lambda$~: en effet, puisque $f_\lambda$ est
quasi-séparé, l'immersion diagonale $X_\lambda\to X'_\lambda$ est
quasi-compacte, et localement de présentation finie puisque $f_\lambda$ est
de type fini (I, 4.3 et I, 5.4, (iii))~; il résulte alors de (1.8.4.1) que
$\Delta_\lambda$ est localement constructible, donc constructible puisque
$X'_\lambda$ est quasi-compact. On peut maintenant appliquer (8.3.12) en
remplaçant $S_\lambda$ et $Z_\lambda$ par $X'_\lambda$ et $\Delta_\lambda$
respectivement.
\end{proof}

\begin{theorem}[8.10.5]
\EGAReaderTarget{ega4.c04.thm.8-10-5}{IV.8.10.5-fr}
Supposons $S_0$ quasi-compact, $X_\alpha$ et $Y_\alpha$ de présentation finie
sur $S_\alpha$, et soit $f_\alpha:X_\alpha\to Y_\alpha$ un
$S_\alpha$-morphisme. Considérons, pour un morphisme, la propriété d'être~:
\begin{enumerate}
\item[(i)] un isomorphisme~;
\item[(i \emph{bis})] un monomorphisme~;
\item[(ii)] une immersion~;
\item[(iii)] une immersion ouverte~;
\item[(iv)] une immersion fermée~;
\item[(v)] séparé~;
\item[(vi)] surjectif~;
